\section{边界：这篇论文没有做什么}\label{sec:boundary}

\subsection{这一章要解决什么}
读完前十一讲，最容易犯的错是"把纲领当定理"。本章把原文的留白逐条列出来，
再把后世补齐的部分做成时间线，最后给出五个仍然开放的问题与三遍读法。
判断标准很简单：\textbf{凡是原文没有给出构造、只给了一句断言的，都要降级。}

\subsection{原文自己的留白（五处）}
\begin{enumerate}[label=B\arabic*,leftmargin=*]
  \item \textbf{函子性/缝合}。p.397 只有一句话："It is obvious at this point that…
        the conformal field theory… is the same as the one obtained by quantizing the
        Chern--Simons theory." 没有给缝合（gluing）的论证，也没有给映射类群作用的完备性证明 \src{397}。
  \item \textbf{一般手术下的标架追踪}。p.390--391 只对简单情形验证了标架反常；
        一般手术链环的 $T^nST^m$ 追踪留待另文 \src{390--391}。
  \item \textbf{半形式修正}。全文没有引入 $K^{1/2}$。$k\to k+h^\vee$ 只从
        (2.16)(2.17) 的 $\eta$ 修正与 (4.21) 的 $q$ 间接出现；机制是后世补的 \src{359,382}。
  \item \textbf{量子群}。p.372 说经典解会得到"量子群"的某种东西，用 \emph{should}，
        并明说 "deferring a fuller treatment"，没有构造 \src{372}。
  \item \textbf{一般群与一般表示}。大部分显式计算在 $SU(N)$ 定义表示下做；
        一般单群只给原则 \src{355}。
\end{enumerate}

\subsection{后世补齐（时间线）}
\begin{center}
\begin{tabular}{@{}llp{8.4cm}@{}}
\toprule
年份 & 工作 & 补了什么 \\
\midrule
1989 & Moore--Seiberg & 用手征代数给 Verlinde 公式的物理证明 \\
1989 & Tsuchiya--Ueno--Yamada & 共形块丛的数学构造，Verlinde 的数学路线之一 \\
1991 & Reshetikhin--Turaev & 用表示论/量子群构造三维不变量（组合版本） \\
1994 & Faltings & Verlinde 公式的代数几何证明（rank-level 对偶、GIT 正性） \\
2000 & Khovanov & Jones 多项式的范畴化 \\
2007 & Kapustin--Witten & 四维扭曲理论框架（几何 Langlands） \\
2011 & Witten & 用五维/四维构造解释 Khovanov 同调 \\
\bottomrule
\end{tabular}
\end{center}

\begin{center}
\begin{tabular}{@{}p{4.2cm}p{3.4cm}p{6.0cm}@{}}
\toprule
原文留白 & 后世补法 & 本讲义对应 \\
\midrule
函子性/缝合（p.397） & Segal 模函子（1988 已提出） & 第 \ref{sec:blocks} 讲 \S5.5 \\
标架追踪（p.390） & RT 1991、Turaev 1992 & 第 \ref{sec:surgery} 讲 \S8.4 \\
半形式（全文无） & 几何量子化文献 & 第 \ref{sec:prequant} 讲 \S4.5 \\
量子群（p.372） & Drinfeld--Jimbo & 第 \ref{sec:wilson} 讲 \S6.4 \\
一般群（p.355） & 后续文献 & 第 \ref{sec:cs} 讲 \S2.6 \\
\bottomrule
\end{tabular}
\end{center}

\subsection{五个仍然开放的问题}
\begin{enumerate}[leftmargin=*]
  \item \textbf{TUY 1989 证到什么程度？} 是证明了公式，还是只在 A 型给出共形块丛
        与 Hitchin/KZ 联络下的维数？
  \item \textbf{$k\to k+h^\vee$ 的符号与来源}：S 矩阵分母的 $k+h^\vee$ 与半形式修正
        $K^{1/2}=\cL_{-\rho}$ 是否同一个 $k$？两条路线给的 $k$ 是否对齐？
  \item \textbf{可积截断与量子群}：$2j\le k$ 的截断（$N_{111}=0$）如何与 $U_q(\mathfrak{sl}_2)$
        的表示范畴精确对应？本文只有一句话。
  \item \textbf{函数性}：共形块丛的射影平坦联络能否升级为严格函子（Segal 模函子）？
        p.397 的 "obvious" 能否补成证明？
  \item \textbf{与全息的关系}：本文的 CS 理论与 AdS$_3$/CFT$_2$ 有没有比"3d 引力 = CS"
        更深的联系？（第 \ref{sec:duality} 讲的结论：没有，本文不是全息。）
\end{enumerate}

\subsection{三遍读法}
\begin{enumerate}[leftmargin=*]
  \item \textbf{第一遍（结构）}：只读第 \ref{sec:intro}、\ref{sec:cs}、\ref{sec:verlinde} 讲，
        抓住"作用量 $\to$ 相空间 $\to$ 共形块 $\to$ 融合规则"四步。
  \item \textbf{第二遍（推导）}：读第 \ref{sec:phase}--\ref{sec:surgery} 讲，动笔复算
        (3.1)(3.2)(4.28)(4.51)，做每章练习；卡住处查附录 \ref{app:exercises}。
  \item \textbf{第三遍（核对）}：回到具体公式核对约定（$k$ 的归一化、$\mathrm{Tr}$ 的表示、
        半形式修正），并区分"原文说 / 我们补 / 后世"；约定见附录 \ref{app:conventions}，
        已核对的事实与译本笔误见附录 \ref{app:errata}。
\end{enumerate}

\subsection{一句总结}
这篇论文给的是"一个可精确求解的构造 + 一整套自洽计算 + 若干断言"。
把它读成"已经严格的定理集"会处处别扭；把它读成"物理学家给出的、后来被数学严格化的纲领"，
每一节的位置就都清楚了。本讲义的工作，就是把这根骨架两端的缝——物理直觉与数学细节——补上。

\begin{boundary}
\begin{itemize}
  \item 已证（原文或引用）：构造与自洽性；若干例子上的手术公式；$SU(2)$ 的低 $k$ 情形。
  \item 我们补：D-x 推导、约定换算（附录 \ref{app:conventions}）、十步链（第 \ref{sec:verlinde} 讲）。
  \item 仍开放：上节五个问题；以及"射影平坦丛的严格函子性"。
\end{itemize}
\end{boundary}
